(1) 球面对称扰动满足 \(\partial_t^2p=c^2r^{-2}\partial_r(r^2\partial_rp)\)。将 \(p=(P/r)e^{i(kr-\omega t)}\) 代入可得 \(\partial_t^2p=-\omega^2p\)、右侧为 \(-c^2k^2p\)，因此 \(k=\omega/c\) 时恰为向外传播的球面解。
(2) 线性欧拉方程为 \(\rho\partial_tu_r=-\partial_rp\)。对 \(e^{-i\omega t}\) 约定，\(\tilde u_r=(1/(i\omega\rho))\partial_r\tilde p=(\tilde p/(\rho c))(1+i/(kr))\)。第一项与声压同相，第二项相差 \(90^\circ\)，并在近声源区域相对增大。
(3) 时间平均声强 \(\langle I_r\rangle=\frac12\Re[\tilde p\,\tilde u_r^*]=|P|^2/(2\rho c r^2)\)，其中虚部不贡献周期平均净流。乘以球面面积 \(4\pi r^2\)，得到 \(P_{\rm acoustic}=2\pi |P|^2/(\rho c)\)，与半径无关，体现无吸收空间的声能守恒。
(4) \(kr\gg1\) 时 \(\tilde u_r\approx\tilde p/(\rho c)\)，局部近似平面行波；\(kr\ll1\) 时 \(i/(kr)\) 项占优，局部速度几乎与声压正交，产生来回交换的反应能量。净向外通量仍由同相部分决定，不可用速度总幅值直接乘压强幅值当作平均辐射功率。
